% !TEX root = ../Main.tex

% === TexCount-Definitions ===
% (Ignore Headings)
%TC:macro \chapter [ignore]
%TC:macro \section [ignore]
%TC:macro \subsection [ignore]
%TC:macro \subsubsection [ignore]
%TC:macro \includesvg [ignore]
% (Ignore Figures and Tables)
%TC:envir figure [ignore] ignore
%TC:envir table [ignore] ignore

% === Content ===
\newpage
% !TEX root = ../../Main.tex
% ==============================================================================
% Kapitel 1: Einleitung
% ==============================================================================
% Description: Einleitung zur Formatierung der Seminararbeit
% ==============================================================================

\chapter{Einleitung}
\label{chap:einleitung}

Das hier vorliegende Dokument soll als Vorlage für die Gestaltung der im Seminarkurs \emph{Ausgewählte Themen der Medizinisch-Technischen Informatik} erstellten schriftlichen Ausarbeitungen dienen.

Es ist zu beachten, dass das Formatieren des Textes als einer der letzten Schritte der Ausarbeitung durchgeführt wird, da dieser Schritt erfahrungsgemäß viel Zeit in Anspruch nimmt und daher nur einmalig ausgeführt werden sollte.

Der vorliegende Entwurf wurde mit \LaTeX{} unter Verwendung der HSRTReport-Dokumentklasse erstellt, welche auf der KOMA-Script-Klasse \texttt{scrreprt} basiert. Diese Vorlage bietet eine professionelle und konsistente Formatierung für wissenschaftliche Arbeiten.

Alle in diesem Dokument enthaltenen Hinweise zur Gestaltung des Dokumentes dienen als Referenz für die korrekte Verwendung des Templates. Die Vorlage selbst übernimmt bereits die meisten Formatierungsanforderungen automatisch.

% ==============================================================================
% End of Kapitel 1
% ==============================================================================

%!TEX root = ../../main.tex
% ==============================================================================
% CHAPTER 2: WISSENSCHAFTLICHE EVIDENZ
% ==============================================================================
\chapter{Wissenschaftliche Evidenz für Enhancement durch Hirnstimulation}
\label{chap:evidenz}

% LITERATURVORSCHLÄGE FÜR KAPITEL:
% - Simonsmeier et al. (2018): "Electrical brain stimulation improves learning" [Meta-Analyse]
% - Reis et al. (2009): "Noninvasive cortical stimulation enhances motor skill acquisition" [motorisch, 1700+ Zit.]
% - Akkad et al. (2021): "Increasing human motor skill acquisition by driving theta-gamma" [innovative Studie]
% - Cavaleiro et al. (2020): "Memory and Cognition-Related Neuroplasticity" [Mechanismen]
% - Frontiers meta-analysis (2024): TMS effects on cognition
% - Vergallito et al. (2022): "Inter-Individual Variability in tDCS Effects" [Variabilität, 154 Zit.]

\section{Neurobiologische Wirkmechanismen}
\label{sec:wirkmechanismen}

% LITERATUR: Cavaleiro 2020, Woods 2016 technical guide, Niksche & Paulus reviews

% \subsection{Grundprinzipien der transkraniellen Hirnstimulation}
% \label{subsec:grundprinzipien}
% - Elektromagnetische Induktion
% - Modulation kortikaler Erregbarkeit
% - Unterschiede tDCS (konstant), tACS (oszillatorisch), TMS (magnetisch)
% - Relevante Hirnregionen (DLPFC, M1, parietal cortex)

% Subsubsection 2.1.1.1: tDCS Mechanismus
% \subsubsection{Transkranielle Gleichstromstimulation (tDCS)}
% \label{subsubsec:tdcs_mech}
% - Anodale vs. kathodale Stimulation
% - Membranpotenzial-Verschiebung
% - Stromfluss und elektrische Feldverteilung
% - Parameter: Intensität (1-2 mA), Dauer (20-30 min), Elektroden-Montage

% Subsubsection 2.1.1.2: tACS Mechanismus
% \subsubsection{Transkranielle Wechselstromstimulation (tACS)}
% \label{subsubsec:tacs_mech}
% - Frequenzabhängige Modulation
% - Entrainment neuronaler Oszillationen
% - Theta (4-8 Hz), Gamma (40 Hz), andere Bänder
% - Theta-Gamma-Kopplung für Gedächtnis/Lernen

% Subsubsection 2.1.1.3: TMS/rTMS Mechanismus
% \subsubsection{Transkranielle Magnetstimulation (TMS/rTMS)}
% \label{subsubsec:tms_mech}
% - Faradaysche Induktion
% - Aktionspotenziale in kortikalen Neuronen
% - High-frequency (>5 Hz) exzitatorisch, Low-frequency (<1 Hz) inhibitorisch
% - iTBS vs. cTBS Protokolle

% \subsection{Zelluläre und molekulare Plastizitätsmechanismen}
% \label{subsec:plastizität}

% LITERATUR: Cavaleiro 2020, Frontiers molecular neuroscience 2020, Esser 2006

% Subsubsection 2.1.2.1: Long-Term Potentiation (LTP)
% \subsubsection{Langzeitpotenzierung und -depression}
% \label{subsubsec:ltp_ltd}
% - LTP als synaptische Verstärkung
% - LTD als synaptische Abschwächung
% - iTBS induziert LTP-ähnliche Effekte
% - cTBS induziert LTD-ähnliche Effekte
% - Messung: MEP (Motor Evoked Potentials)

% Subsubsection 2.1.2.2: Neurotrophische Faktoren
% \subsubsection{Neurotrophische Faktoren und Genexpression}
% \label{subsubsec:neurotrophisch}
% - BDNF (Brain-Derived Neurotrophic Factor)
% - c-fos, zif268 (immediate early genes)
% - NMDA-Rezeptoren
% - Hochregulation durch Stimulation (besonders iTBS)

% Subsubsection 2.1.2.3: Transmitter-Modulation
% \subsubsection{Neurotransmitter-Modulation}
% \label{subsubsec:transmitter}
% - GABAerge und glutamaterge Systeme
% - Dopamin und Lernverhalten (Reward-based Learning)
% - Monoamine (Serotonin, Noradrenalin)
% - Inter-individuelle Variabilität durch Transmitter-Polymorphismen

% \subsection{Zeitliche Aspekte: Phasen des Lernens}
% \label{subsec:zeitliche_aspekte}

% LITERATUR: Reis 2009, Debarnot 2019, Zimerman 2013

% Subsubsection 2.1.3.1: Enkodierung und Online-Lernen
% \subsubsection{Enkodierung und Online-Lernen}
% \label{subsubsec:enkodierung}
% - Effekte während aktiver Aufgabe
% - Stimulation während Trainingsphase
% - Hemmt oder fördert Online-Lernen je nach Timing/Polarität
% - Studien: [Refs: Reis 2009 - motorisch, Stagg 2011]

% Subsubsection 2.1.3.2: Konsolidierung und Offline-Effekte
% \subsubsection{Konsolidierung und Offline-Effekte}
\label{subsubsec:konsolidierung}
% - Verstärkung zwischen Trainings-Sessions
% - Sleep-dependent consolidation
% - tDCS während Tiefschlaf: Slow Oscillations induzieren
% - Externe Effekte: 1h, 24h, 1 Woche post-stimulation
% - Studien: [Reis 2009, Marshall & Born 2007]

\section{Empirische Befunde nach kognitiver Domäne}
\label{sec:befunde_domänen}

% LITERATURVORSCHLÄGE: Simonsmeier 2018, Reis 2009, Akkad 2021, Lopez-Alonso 2025

% \subsection{Motorisches Lernen und Fertigkeitserwerb}
\label{subsec:motorisch}

% LITERATUR: Reis 2009 (1700 Zit.), Akkad 2021, Maceira-Elvira 2024, Lopez-Alonso 2025

% Subsubsection 2.2.1.1: Anodales tDCS über M1
% \subsubsection{Anodale tDCS über motorischem Kortex (M1)}
\label{subsubsec:motor_tdcs}
% - Effekte auf Konsolidierung (nicht Online-Lernen)
% - Protokoll: 2 mA, 20 min, posttraining
% - Effektgrößen: Cohen's d = 0.5-1.0
% - Persistenz: bis zu 24h
% - Alter-abhängig: auch ältere Probanden profitieren

% Subsubsection 2.2.1.2: Theta-Gamma-tACS über M1/sensomotorischer Kortex
% \subsubsection{Theta-Gamma-tACS für motorisches Lernen}
\label{subsubsec:motor_tacs}
% - Innovative Frequenzkopplung
% - Phase-Amplitude-Coupling (4-8 Hz Theta moduliert 40 Hz Gamma)
% - Replikationsstudie [Akkad 2021]
% - Effekte: Beschleunigung, Persistenz über 1h
% - Mechanismus: Theta-Gamma relevant für hippokampales Lernen

% Subsubsection 2.2.1.3: Komplexe motorische Skills
% \subsubsection{Komplexe motorische Skills (negative Befunde)}
\label{subsubsec:komplexe_motor}
% - Golf-Putting Studie [Lopez-Alonso 2025]
% - tDCS zeigte KEINE zusätzliche Verbesserung über Training
% - Implikation: Einfache vs. komplexe Skills unterschiedlich
% - Limitationen: Realwelt-Komplexität, Multi-Link Bewegungen

% \subsection{Akademische Fähigkeiten (Mathematik, Sprache)}
\label{subsec:akademisch}

% LITERATUR: Simonsmeier et al. 2018 (133 Zit., Meta-Analyse), Cohen Kadosh Arbeiten, Sarkar 2014

% Subsubsection 2.2.2.1: Timing der Stimulation
% \subsubsection{Timing: Stimulation während Lernphase vs. Test}
\label{subsubsec:timing_learning}
% - Simonsmeier Meta-Analyse [2018]: Stimulation während Lernen > während Test
% - Mechanismus: LTP-Induktion während aktiver Verarbeitung
% - Effektgrößen: SMD = 0.3-0.7 für akademische Fähigkeiten

% Subsubsection 2.2.2.2: Mathematische Kompetenzen
% \subsubsection{Mathematische Fähigkeiten}
\label{subsubsec:mathe}
% - tRNS (Random Noise) über "Mathezentrum" (DLPFC/IPS)
% - Schmerzfreie Alternative zu tDCS
% - Verbesserung Reaktionszeit + Genauigkeit
% - [Cohen Kadosh Arbeiten, erwähnt in Spektrum.de]

% Subsubsection 2.2.2.3: Sprachfähigkeiten
% \subsubsection{Sprachfähigkeiten}
\label{subsubsec:sprache}
% - Weniger Forschung als Mathematik
% - tDCS über Sprach-Arealen (Broca, Wernicke)
% - Stimulation während Vokabel-Lernen
% - Effektgrößen: moderate bis klein

% \subsection{Arbeitsgedächtnis und exekutive Funktionen}
\label{subsec:workingmem}

% LITERATUR: Hoy et al. 2016, Senkowski & Sobirey 2022, Imburgio & Orr 2018

% Subsubsection 2.2.3.1: tDCS vs. tACS Vergleich
% \subsubsection{Vergleich tDCS vs. tACS auf Arbeitsgedächtnis}
\label{subsubsec:tdcs_vs_tacs}
% - Systematische Review [Senkowski 2022]: 43 Studien, 1826 Teilnehmer
% - Single-Session tDCS: kaum signifikant
% - Multi-Session tDCS: moderate Effekte (SMD ~0.3)
% - tACS (v.a. Gamma): robustere Effekte (SMD ~0.4-0.6)
% - Frequenz-Abhängigkeit: Gamma (40 Hz) > andere

% Subsubsection 2.2.3.2: iTBS auf DLPFC
% \subsubsection{Intermittierende Theta-Burst-Stimulation (iTBS) auf DLPFC}
\label{subsubsec:itbs_dlpfc}
% - Hoy et al. (2016): iTBS > tDCS auf 3-back Working Memory Task
% - Theta & Gamma Synchronisation erhöht
% - Effektgrößen größer als tDCS
% - Schnellere Behandlung (10min statt 20min)

% Subsubsection 2.2.3.3: Exekutive Funktionen
% \subsubsection{Exekutive Funktionen (Inhibition, Flexibilität)}
\label{subsubsec:exek_funkt}
% - tDCS Effekte: limitiert auf Refresh (WM)
% - Inhibition: kleine bis keine Effekte [Imburgio & Orr 2018]
% - Kognitive Flexibilität: wenig untersucht
% - Kontext-abhängigkeit der Effekte

% \subsection{Praktische Anwendungen in der Realwelt}
\label{subsec:realwelt_apps}

% LITERATUR: Clark et al. 2012, Lopez-Alonso 2025, Violante 2024

% Subsubsection 2.2.4.1: Threat-Detection-Training (Militär)
% \subsubsection{Threat-Detection-Training}
\label{subsubsec:threat_detection}
% - Clark et al. (2012): Virtual Reality + tDCS
% - 96 Probanden, Erkennung versteckter Objekte
% - Signifikante Verbesserung + Trainingstransfer
% - DARPA finanzierte Forschung
% - Praktische Implikation: funktioniert in komplexer Umgebung

% Subsubsection 2.2.4.2: Negative Befunde und Realwelt-Limitationen
% \subsubsection{Negative Befunde und Realwelt-Limitationen}
\label{subsubsec:negative_realwelt}
% - Lopez-Alonso (2025): Golf-Putting keine zusätzliche Verbesserung
% - Violante et al. (2024): "modest evidence" für Neuro-Enhancement
% - Komplexe Skills brauchen Integration verschiedener Systeme
% - Transfer-Problem: Lab ≠ Realwelt

\section{Kritische Limitationen der Evidenz}
\label{sec:limitationen}

% LITERATURVORSCHLÄGE: Vergallito 2022, Chew 2015, Maceira-Elvira 2024, Violante 2024

% \subsection{Inter-individuelle Variabilität in Stimulationsresponse}
\label{subsec:variabilität}

% LITERATUR: Vergallito et al. 2022 (154 Zit.), Chew et al. 2015 (317 Zit.), Filmer et al. 2014

% Subsubsection 2.3.1.1: Anatomische Unterschiede
% \subsubsection{Anatomische Unterschiede}
\label{subsubsec:anat_diff}
% - Kortikale Gyrierung (Gyri/Sulci Patterns)
% - Kortexdicke in spezifischen Regionen
% - Faserverbindungen zwischen Regionen
% - Vorhersage: Dicke in mittlerem Stirnfurchen prädiziert tDCS-Response [Filmer]
% - Modellierung: FEM (Finite-Elemente) kann Feldverteilung individualisieren

% Subsubsection 2.3.1.2: Genetische und biochemische Faktoren
% \subsubsection{Genetische und biochemische Faktoren}
\label{subsubsec:genetik}
% - BDNF val66met Polymorphismus
% - Transmitter-Polymorphismen (COMT, DAT)
% - Baseline-GABA/Glutamat Ratio
% - Limitierte Forschung, aber vielversprechend für Personalisierung

% Subsubsection 2.3.1.3: Funktionelle Baseline-Unterschiede
% \subsubsection{Funktionelle Baseline-Unterschiede}
\label{subsubsec:baseline_func}
% - Resting-State Konnektivität
% - Theta-Power im Frontallappen (prädiziert rTMS-Response bei Depression)
% - Baseline-Leistung auf kognitiver Aufgabe

% \subsection{Baseline-Leistung und Lerner-Typ}
\label{subsec:baseline_leistung}

% LITERATUR: Maceira-Elvira 2024, Sarkar 2014

% Subsubsection 2.3.2.1: Wer profitiert am meisten?
% \subsubsection{Baseline-abhängige Responsiveness}
\label{subsubsec:baseline_response}
% - Maceira-Elvira (2024): "Native learning ability, not age, determines response"
% - Personen mit niedrigen Ausgangsfähigkeiten profitieren MEHR
% - Hochleister: KEINE zusätzliche Verbesserung durch tDCS
% - Implikation: Enhancement-Effekt abhängig vom Starting Point

% Subsubsection 2.3.2.2: Math-Anxiety Beispiel
% \subsubsection{Kontext-abhängige Effekte: Math-Anxiety}
\label{subsubsec:math_anxiety}
% - Sarkar et al. (2014): tDCS hilft bei Mathe-Angst
%   - Verbesserung Reaktionszeit + reduziertes Cortisol
% - Aber: Bei Nicht-Angst-Personen verschlechtert sich Leistung!
% - Implikation: Effekte nicht universal, stark kontext-abhängig

% \subsection{Nur 39-45\% sind Responder}
\label{subsec:responder_rate}

% LITERATUR: Vergallito 2022, Chew 2015, Filmer 2014

% Subsubsection 2.3.3.1: Empirische Evidenz
% \subsubsection{Responder vs. Non-Responder Phänomenologie}
\label{subsubsec:responder_phenom}
% - Vergallito (2022): Nur 39-45% zeigen erwartete Effekte
% - Cluster-Analysen zeigen bimodale Verteilungen
% - Chew (2015): Intra-individual Variability ICC = -0.50 (negligible!)
% - Problem: Nicht zuverlässig vorhersagbar mit Current Protocols

% Subsubsection 2.3.3.2: Implikationen für praktische Anwendung
% \subsubsection{Implikationen für praktische Anwendung}
\label{subsubsec:responder_impl}
% - Screening notwendig zur Identification von Respondern
% - Adaptives Protokoll-Adjustments
% - Oder: Personalisierung VOR Behandlung (individuelles FEM-Modell, fMRI)

% \subsection{Transferdefizite und Aufgabenspezifität}
\label{subsec:transfer}

% LITERATUR: Violante 2024, Lopez-Alonso 2025

% Subsubsection 2.3.4.1: Begrenzte Generalisierbarkeit
% \subsubsection{Begrenzte Generalisierbarkeit der Effekte}
\label{subsubsec:generalisierbar}
% - Effekte meist aufgabenspezifisch
% - Transfer zu anderen Aufgaben limitiert
% - Transfer zu anderen Domänen: unklar
% - Implikation: Nicht "allgemeine Intelligenz-Verbesserung"

% Subsubsection 2.3.4.2: Realwelt-Transfer
% \subsubsection{Transfer von Lab zu Realwelt}
\label{subsubsec:lab_realwelt}
% - Threat-Detection erfolgreich [Clark 2012]
% - Golf-Putting KEIN Transfer [Lopez-Alonso 2025]
% - Komplexität des Realwelt-Skills kritisch
% - Context-Dependence der Stimulation-Effekte

% \subsection{Fehlende Langzeitstudien und Persistenz}
\label{subsec:langzeit}

% LITERATUR: Brühl 2019, Woodham 2024 (nur 10 Wochen), Violante 2024

% Subsubsection 2.3.5.1: Kurzfristige vs. Langfristige Effekte
% \subsubsection{Persistenz von Effekten}
\label{subsubsec:persistenz}
% - Kurzfristig: Gut dokumentiert (bis 24h post-stimulation)
% - Mittelfristig (1-4 Wochen): Noch nicht vollständig untersucht
% - Langfristig (Monate/Jahre): Kaum Daten
% - Frage: Was passiert nach Beendigung der Stimulation?

% Subsubsection 2.3.5.2: Lifetime Safety bei Gesunden
% \subsubsection{Lifetime Safety bei wiederholter Anwendung bei Gesunden}
\label{subsubsec:lifetime_safety}
% - Kurzfristige Nebenwirkungen: gut dokumentiert, mild
% - Langfristige Effekte unbekannt [Brühl 2019]
% - Besondere Bedenken: Hirnentwicklung (<25 Jahre)
% - Frage: Könnte wiederholte Stimulation Neuroplastizität negativ beeinflussen?
% - Forschungsbedarf: Prospektive Langzeitstudien

%!TEX root = ../../main.tex
% ==============================================================================
% CHAPTER 3: PRAKTIKABILITÄT
% ==============================================================================
\chapter{Praktikabilität in Bildungs- und Trainingskontexten}
\label{chap:praktikabilität}

% LITERATURVORSCHLÄGE FÜR KAPITEL:
% - Cappon et al. (2023): "Home-based tES training program" [Machbarkeit, 10 Zit.]
% - Woodham et al. (2024): "Home-based tDCS in MDD" [Nature Medicine, 62 Zit., praktisches Modell]
% - Woods et al. (2015/2016): "Technical guide to tDCS" [Standardwerk, 1700+ Zit.]
% - Thair et al. (2017): "Transcranial Direct Current Stimulation (tDCS)" [Praktische Details]
% - Caulfield et al. (2022): "Optimized electrode positions" [Technische Optimierung]

\section{Technische Aspekte und Durchführbarkeit}
\label{sec:tech_durchführbarkeit}

% LITERATUR: Woods 2015/2016, Thair 2017, Caulfield 2022

% \subsection{Kosteneffizienz und ökonomische Zugänglichkeit}
\label{subsec:kosten}

% LITERATUR: neurocare academy, Woodham 2024, Cappon 2023

% Subsubsection 3.1.1.1: Gerätekosten
% \subsubsection{Gerätekosten und Verfügbarkeit}
\label{subsubsec:gerätekosten}
% - tDCS Consumer-Geräte: $200-500
% - rTMS klinische Geräte: $10,000-50,000
% - tACS/tRNS: mittleres Preissegment
% - Vergleich:
%   - Private Nachhilfe/Einzelunterricht: €30-60/h
%   - Elite-Trainingsprogramme: €1000+/Monat
%   - tDCS: ~€5-10 pro Sitzung (Amortisation)

% Subsubsection 3.1.1.2: Kostenvergleich zu Alternativen
% \subsubsection{Kostenvergleich zu alternativen Interventionen}
\label{subsubsec:kostenvergleich}
% - Vs. Einzelunterricht: tDCS kostengünstiger bei Skalierung
% - Vs. Medikamente (Modafinil, etc.): vergleichbar
% - Vs. kognitive Trainings-Apps: tDCS höher initial, aber kombinierbar

% \subsection{Heimbasierte und dezentrale Anwendung}
\label{subsec:heimbasiert}

% LITERATUR: Cappon 2023 (10 Zit., Pilotprogramm), Woodham 2024 (Nature Med, 62 Zit., 174 Probanden)

% Subsubsection 3.1.2.1: Erfolgreiche Pilotprojekte
% \subsubsection{Erfolgreiche Pilotprojekte mit Heimanwendung}
\label{subsubsec:pilotprojekte}
% - Cappon et al. (2023): Home-based tES Trainings-Programm
%   - Laien-Administration nach Trainig
%   - Adhärenz: 98-100%
%   - Sicherheit demonstriert
%   - Downloadbare Materialien verfügbar
% - Woodham et al. (2024, Nature Medicine): Home-based tDCS bei MDD
%   - 174 Teilnehmer, remote supervision
%   - 10 Wochen Protokoll
%   - Hohe Akzeptanz, gute Effektivität

% Subsubsection 3.1.2.2: Remote Supervision und Monitoring
% \subsubsection{Remote Supervision und Monitoring}
\label{subsubsec:remote_supervision}
% - Möglichkeit mit modernen Technologien
% - Telemedizin-Modelle existieren
% - Dokumentation via Apps/Cloud
% - Sicherheit: Contact person im Notfall nötig

% Subsubsection 3.1.2.3: Skalierbarkeit für breite Populationen
% \subsubsection{Skalierbarkeit für breite Populationen}
\label{subsubsec:skalierbarkeit}
% - Implication: Theoretisch skalierbar
% - Unterschied:
%   - Élite-Training (kleine Gruppen, hohe Personalkosten): möglich
%   - Mainstream-Bildung (Millionen Schüler): logistisch schwierig
%   - Hybrid-Modell: vielversprechend (Schulungen + Home-Use)

% \subsection{Sicherheit und Nebenwirkungsprofil}
\label{subsec:sicherheit}

% LITERATUR: Woods 2016 (technischer Standard), Wassermann-Kriterien, McCambridge et al. Sicherheit-Reviews

% Subsubsection 3.1.3.1: Häufige milde Nebenwirkungen
% \subsubsection{Häufige und milde Nebenwirkungen}
\label{subsubsec:mild_nebenwirkungen}
% - Kopfschmerzen: 10-20% (mit rezeptfreien Schmerzmitteln behandelbar)
% - Kopfhaut-Irritation/Brennen: lokal, vorübergehend
% - Muskelzuckungen (bei TMS): während Stimulation, gewöhnung nach wenigen Sessions
% - Müdigkeit/Benommenheit: selten, vorübergehend (<1h)
% - Übelkeit: sehr selten

% Subsubsection 3.1.3.2: Schwerwiegende Komplikationen
% \subsubsection{Schwerwiegende Komplikationen}
\label{subsubsec:severe_complications}
% - Epileptische Anfälle: extrem selten (<0.1% bei compliance mit Wassermann-Kriterien)
% - Keine bekannten Fälle permanenter Hirnschädigung
% - Vergleich zu Medikamenten: günstiger Effekt-Nebenwirkung-Profile
% - Kontraindikationen gut etabliert

% Subsubsection 3.1.3.3: Vergleich zu pharmakologischen Enhancern
% \subsubsection{Vergleich zu pharmakologischen Enhancern}
\label{subsubsec:pharm_vs_stim}
% - tDCS: keine systemischen Effekte (lokal wirksam)
% - Modafinil/Amphetamine: systemische Nebeneffekte (kardiovaskular, GI, etc.)
% - Koffein: Entzugssymptome, Toleranz
% - Fazit: tDCS sicherere Alternative aus Nebenwirkungs-Perspektive

\section{Praktische Implementierungshürden}
\label{sec:hürden}

% LITERATURVORSCHLÄGE: Vergallito 2022, Woods 2016, Meinzer 2024, Cappon 2023

% \subsection{Personalisierung und Vorhersage von Responsiveness}
\label{subsec:personalisierung}

% LITERATUR: Vergallito 2022, Meinzer 2024 (neuroimaging context), computational modeling papers

% Subsubsection 3.2.1.1: Das Problem: 60\% Non-Responder
% \subsubsection{Das Problem der 60\% Non-Responder}
\label{subsubsec:nonresponder_problem}
% - Vergallito (2022): Nur 39-45% zeigen erwartete Effekte
% - Current Standardprotokolle: "One-Size-Fits-All" funktioniert nicht
% - Screening nötig ODER Individualisierung VOR Behandlung

% Subsubsection 3.2.1.2: Anatomische Personalisierung
% \subsubsection{Anatomische Personalisierung}
\label{subsubsec:anat_personalisierung}
% - MRT-basierte Modelle: Individuelle Elektrodenplatzierung
% - Finite-Elemente-Modelle (FEM): Stromfluss berechnen
% - Toolboxen: SimNIBS, COMETS (MATLAB)
% - Kosten/Nutzen: Trade-off zwischen Präzision und praktische Durchführbarkeit

% Subsubsection 3.2.1.3: Funktionelle und genetische Personalisierung
% \subsubsection{Funktionelle und genetische Biomarker}
\label{subsubsec:biomarker_personal}
% - Theta-Power im Frontallappen (prädiziert rTMS-Response)
% - BDNF-Polymorphismen, COMT, DAT
% - EEG-basierte adaptive Protokolle (Closed-Loop)
% - Forschungsbedarf: Validierung dieser Biomarker

% \subsection{Standardisierung und Schulung von Anwendern}
\label{subsec:standardisierung}

% LITERATUR: Cappon 2023 (Trainings-Programm), neurocare academy, IFCN Guidelines

% Subsubsection 3.2.2.1: Notwendigkeit von Standardisierten Protokollen
% \subsubsection{Standardisierte Protokolle und Guidelines}
\label{subsubsec:standardisierte_protokolle}
% - Unterschiede aktuelle Praxis: große Variabilität
% - IFCN Guidelines (International Federation of Clinical Neurophysiology)
% - Wassermann-Kriterien für Sicherheit
% - Notwendig: Einigung auf "Best Practice" Protokolle

% Subsubsection 3.2.2.2: Trainingsmaterialien für nicht-medizinisches Personal
% \subsubsection{Trainingsmaterialien für Laien}
\label{subsubsec:training_laien}
% - Cappon (2023) hat Programm entwickelt: Video, Quizze, Checklisten
% - Erfolgreiche Durchführbarkeit demonstriert
% - Unterschied: Laien können Geräte anlegen, aber müssen medizinisch überwacht werden

% Subsubsection 3.2.2.3: Regulatorische Fragen zur Anwender-Qualifikation
% \subsubsection{Regulatorische Fragen}
\label{subsubsec:regulatory_qual}
% - Wer darf tDCS/TMS anwenden?
% - Ärzte, Therapeuten, Trainer, Laien?
% - Versicherungsdeckung, Haftung
% - Zertifizierungspfade nötig

% \subsection{Realistische Effektgrößen und Erwartungsmanagement}
\label{subsec:effektgrößen}

% LITERATUR: Reis 2009, Akkad 2021, Simonsmeier 2018, Violante 2024

% Subsubsection 3.2.3.1: Evidenzbasierte Effektgrößen
% \subsubsection{Empirische Effektgrößen aus Studien}
\label{subsubsec:empirisch_effekt}
% - Motorisches Lernen: Beschleunigung ~10-30% [Reis 2009, Akkad 2021]
% - Akademische Fähigkeiten: SMD = 0.3-0.7 [Simonsmeier 2018]
% - Arbeitsgedächtnis (Multi-Session): SMD ~0.3-0.5
% - NICHT: Keine "Wunder-Effekte"

% Subsubsection 3.2.3.2: Vergleich zu anderen Lerntechniken
% \subsubsection{Benchmarking gegen etablierte Methoden}
\label{subsubsec:benchmarking}
% - Spaced Repetition: SMD ~0.5-1.0
% - Elaboration: SMD ~0.7
% - tDCS: SMD ~0.3-0.7 (abhängig von Domäne)
% - Fazit: Vergleichbar zu anderen Techniken, nicht überlegen

% Subsubsection 3.2.3.3: "Nicht-Ersatz für gutes Lehren"
% \subsubsection{tDCS ersetzt nicht gutes Unterrichten}
\label{subsubsec:not_replacement}
% - tDCS ist ZUSATZ zu optimiertem Unterricht
% - Basissystem muss gut sein (sonst "garbage in, garbage out")
% - Kombination wichtig: tDCS + guter Trainer/Lehrer + strukturiertes Programm

\section{Praktikabilitäts-Fazit: Realistisches Szenario}
\label{sec:praktik_fazit}

% LITERATUR: Violante 2024 (kritikal), Cappon/Woodham für Machbarkeit

% \subsection{Machbarkeit nach Setting}
\label{subsec:machbarkeit_setting}

% Subsubsection 3.3.1.1: Élite-Training und spezialisierte Kontexte
% \subsubsection{Élite-Training (Sport, Musik, militärisch)}
\label{subsubsec:elite_training}
% - Hochmotivierte Populationen, kleinere Gruppen
% - Individuelle Betreuung möglich
% - Kosteneffizienz ggf. gegeben
% - Realistisch: 5-10 Jahre bis etablierte Praxis

% Subsubsection 3.3.1.2: Mainstream-Bildung (Schulen, Universitäten)
% \subsubsection{Mainstream-Bildung}
\label{subsubsec:mainstream_bildung}
% - Massive Skala (Millionen), Kostendruck
% - Heterogene Populationen (nicht alle profitieren)
% - Logistische Herausforderungen
% - Realistisch: 15-25 Jahre oder später für breite Integration

% Subsubsection 3.3.1.3: Home-Use für individuelle Lerner
% \subsubsection{Home-Use für Selbstoptimierung}
\label{subsubsec:homeuse_self}
% - Consumer-Geräte bereits verfügbar
% - Größtes Risiko: unsachgemäße Anwendung ohne Supervision
% - Dual-Use Problem: Medizinisch/Enhancement-Nutzung vermischt

%!TEX root = ../../main.tex
% ==============================================================================
% CHAPTER 5: ZUSAMMENFASSUNG UND FAZIT
% ==============================================================================
\chapter{Zusammenfassung und Fazit}
\label{chap:fazit}

% LITERATURVORSCHLÄGE FÜR KAPITEL:
% - Violante et al. (2024): "Can neurotechnology revolutionize cognitive enhancement?" [überblick, kritisch]
% - Brühl et al. (2019): Balanced Assessment [Sicherheit, Ethik]
% - Cappon/Woodham: praktische Perspektive

\section{Wirksamkeit}

\section{Praktikabilität}

\section{Ethische Implikationen}


%TC:ignore
\iffalse
	\chapter{First Chapter}

	Das \gls{abstract} beschreibt in wenigen Sätzen die Zielsetzung und das Ergebnis der Ausarbeitung. Das Abstract muss sich vollständig auf der Titelseite befinden. Die Zeichensatzformatierung wird in einem eigenen Absatz beschrieben  Das Abstract soll es den Lesern:innen ermöglichen, innerhalb von wenigen Augenblicken zu erfassen, welcher Inhalt hinter der Überschrift steckt und ob das Thema, aus Sicht der Leser:innen, zur weiteren Bearbeitung lohnt. Das Abstract ist keine verbale Beschreibung des Inhaltsverzeichnisses, sondern gibt kurz und knapp z.B. die Zielsetzung (z.B. Hypothese), die eingesetzten Methoden und die erzielten Ergebnisse / Erkenntnisse bekannt. Weitere Hinweise finden Sie außerdem im Vorlesungsskript.

	Beispielverweis auf Quelle \cite{ahrensAbschlussarbeiten}.

	Test-Acronym: \acrshort{gcd}.

	\section{First Section}

	% Beispieltabelle
	\begin{table}[h]
		\centering
		\begin{tabular}{|c|c|}
			\hline
			Spalte 1 & Spalte 2 \\
			\hline
			Inhalt 1 & Inhalt 2 \\
			\hline
		\end{tabular}
		\caption{Beispiel-Tabelle}
		\label{tab:example}
	\end{table}

	\chapter{Second Chapter}

	\section{Another Section}

	% Beispielabbildung
	\begin{figure}[h]
		\centering
		\includegraphics[width=0.5\textwidth]{HSRTReport/Assets/Images/METI.png}
		\caption{MeTI-Logo}
		\label{fig:meti}
	\end{figure}

	% Alles hier wird von TexCount ignoriert.
	%TC:endignore
\fi
